% Source: https://artofproblemsolving.com/community/c3h1712197_interesting_number_theory_problem
% Stable identity: aops-topic-1712197-q01
First simplify the left side: $(ge-fe)\left(\dfrac {g}{e}+\dfrac {f}{e}\right)=e(g-f)\cdot\dfrac{g+f}{e}=g^2-f^2$, an integer. So the condition says
\[
(e+f+g)^2\ \Big|\ g^2-f^2-e^2 .
\]
Let $M=(e+f+g)^2$. Then $0<g^2<M$ and $0<e^2+f^2<M$, so $g^2-(e^2+f^2)$ is the difference of two numbers in the interval $(0,M)$ and has absolute value less than $M$. The only multiple of $M$ in that range is $0$. Hence the condition holds exactly when
\[
e^2+f^2=g^2 ,
\]
i.e.\ when $(e,f,g)$ is a Pythagorean triple with hypotenuse $g$. There are infinitely many such triples, for example $(3k,4k,5k)$ for every positive integer $k$. So there are $\boxed{\text{infinitely many}}$ ordered triples.
